\section{Architecture v3: where authority and progress live}
\label{sec:architecture}

The implementation baseline is the cover revision. CODE statements below name implementation
paths; DOC statements describe requirements or published designs. The test descriptions name
what could exercise an invariant. Unless a dated result is given in Section~\ref{sec:results},
their pass state at this baseline is UNVERIFIED: no product test was run for this edition.

\subsection{Planes and traps}

\doc The product is the \code{kf-*} tree linked into the VMM process. A BAR0 write
records demand through bounded token and ring state. Workers service channel demand; a
register drainer services privileged register operations. BAR1 and BAR2 memory accesses do not
trap. There is no separate isolate process or product IPC plane. Confinement applies to the
whole hypervisor process.\src{architecture}\src{corecode}\src{trapcode}

\attack{No host mapping or blocking operation runs on a vCPU or under a lock it takes.}
{The trap crate's unit tests, core composition tests and source restrictions constrain the
reachable trap path. Hardware trap timing checks a separate latency property.\src{trapcode}}
{A nonblocking-looking call can acquire a lock held by slow host work. A mock host that never
blocks cannot demonstrate the production lock discipline. The old paper records that failure
mode; the restriction must be checked transitively.\src{old}}

\subsection{Hostile input and authored operations}

\doc Guest root and guest userspace are untrusted. Host memory operations are constructed
from admitted ownership, extent, permission and kind information; guest control messages are
not replayed to the host. Host GPU facts come from unprivileged queries or generated open
interfaces. This is the intended authority boundary, with the host driver and GPU trusted.
\src{security}\src{architecture}\src{hostcode}

\attack{No guest-chosen raw address or privilege request creates new host authority.}
{Hostile range/ownership tests, channel privilege tests and the physical-operand hardware
oracle exercise different pieces of this claim.\src{security}\src{physevidence}}
{Checking a decoder cannot establish that every caller uses it. A guest-declared OS identity
also cannot prove a safe mapping scope. Whole-path reachability and multiple-VM fault effects
remain separate review obligations.\src{rulings}\src{security}}

\subsection{Derive, never capture}

\doc The architecture prefers facts read from the host's own unprivileged interface,
computed virtual-device state and rules generated from pinned open headers. A device identity
is not permission to copy another die's answer. Request validation and field-level virtual
support still apply when an answer comes from the host. Published measured exceptions belong
to exact die/driver coverage, not an invisible default table.\src{architecture}\src{constants}

\attack{Every guest-visible field has a source and a defined unsupported case; an unknown die
never silently inherits a neighbour's observation.}
{The constants strategy requires source tags, derivation oracles, wrong-answer mutations and
same-die native acceptance. It explicitly separates the implemented baseline from table
consumers on branches.\src{constants}\src{provenance}}
{A source tag proves provenance only if its source actually supports the field. A boot that
never exercises an optional feature cannot verify that feature, and agreement between a few
dies is not a mechanism for a family rule.\src{constants}}

\subsection{Seven axes, not one compatibility version}

\doc The design separates guest kernel, guest driver, host driver, GPU architecture,
GPU die, VMM and guest OS. An architecture supplies family mechanisms; a die supplies actual
inventory and capability facts. The guest ABI controls decoding; the host ABI controls authored
calls. The VMM supplies memory, interrupts and lifecycle events. These are distinct axes even
when a test holds several fixed.\src{design}\src{constants}\src{driverdesign}

\begin{center}\small
\begin{tabularx}{\textwidth}{@{}lXX@{}}
\toprule Axis & Boundary to inspect & Counterexample worth testing \\
\midrule
Guest kernel & DMA addressing / virtual IOMMU & IOVA mistaken for guest physical address \\
Guest driver & versioned decode / accepted replies & layout changes within a branch \\
Host driver & authored RM ABI & successful syscall with failed in-band status \\
Architecture & class / register / MMU family & different engine or boot protocol \\
Die & runtime facts / measured exceptions & floor-swept or unmeasured die \\
VMM & memory, IRQ, reset adapter & reset origin or wake semantics differ \\
Guest OS & driver protocol and scheduling & Windows registration / recovery path \\
\bottomrule
\end{tabularx}
\end{center}
The counterexamples are proposed attacks, not claims that each has occurred at the baseline.
A matrix that changes one axis at a time is evidence for those cells, not the full Cartesian
product.\src{driverdesign}\src{constants}

\attack{Guest and host ABI choices remain independent, and OS or VMM identity never grants
extra memory authority.}
{The driver matrix tests named host/guest cells; versioned decoder oracles check layouts.
Portability beyond QEMU and x86-64 still needs real adapter evidence.\src{driverdesign}
\src{requirements}}
{A layout may change within a driver branch, and a cross-build cannot exercise mapping,
interrupt or reset semantics. Passing one axis with the others fixed is not the full grid.}

\subsection{Crate map and the frozen predecessor}

\doc The following is a functional map from \code{ARCHITECTURE.md}, not a generated
claim about every dependency edge.\src{architecture}
\begin{center}\small
\begin{tabularx}{\textwidth}{@{}p{42mm}X@{}}
\toprule Crates & Responsibility \\
\midrule
\code{kf-trap}, \code{kf-core} & bounded demand publication and plane composition \\
\code{kf-gsp}, \code{kf-rm} & firmware-facing protocol and resource-manager model \\
\code{kf-chip}, \code{kf-abi} & host facts, family rules and versioned generated layouts \\
\code{kf-host} & one in-process host RM session and authored verbs \\
\code{kf-mem}, \code{kf-cuda} & backing store, VA manager, GPU walk and diff protocol \\
\code{kf-chan} & channel birth, translated work, relay and completions \\
\code{kf-qemu} + C device & QEMU memory, lifecycle, interrupt and device adapter \\
\code{kf-disp}, \code{kf-broker} & emulated display and VMM-neutral broker protocol \\
\code{kf-oprom}, \code{kf-gop-image} & generated boot-display container and build \\
\code{kf-harness} & real-GPU tests without QEMU \\
\code{kf-census}, \code{kf-crec}, \code{kf-trace} & offline tools and historical oracles \\
\code{kf-arch}, \code{kf-linux-raw}, \code{kf-util} & vocabulary, audited OS calls, utilities \\
\bottomrule
\end{tabularx}
\end{center}

\doc The older \code{kayfabe-*} crates remain frozen; the raw-client grader depends on
part of that tree. Product \code{kf-*} crates do not depend on it. Unsafe code is confined to
named adapter crates and \code{*_unsafe.rs} files. A crate inventory is a source-level fact;
it says nothing about whether a feature has a production caller.\src{rules}\src{architecture}

\attack{Product crates retain their dependency and unsafe boundaries; the frozen grader tree
never becomes a hidden product dependency.}
{Workspace manifests and CI boundary checks enforce structural restrictions. Reached-path
oracles and real adapter tests address behaviour separately.\src{rules}\src{architecture}}
{A legal dependency graph can still contain an unused implementation or a production-only
blocking path. Source structure does not establish reachability or hardware correctness.}

\subsection{Channel kinds and the meaning of completion}

\doc Passthrough creates a host twin for user work and lets the real engine execute its
commands. Translated channels rewrite kernel copy work into an owned address window and submit
it to the real engine. Emulated channels execute only work whose meaning belongs entirely to
the VMM. The three kinds differ in authority and completion, not simply speed.\src{kinds}

\attack{No CPU executor substitutes for GPU work; no completion is forged for work submitted
to hardware.}
{The channel suites and real-engine gates check actual execution, content and notification.
The translated-plane design separates address arithmetic from execution.\src{translated}}
{An application can report success without checking its GPU error or result. A harness must
validate the output and completion path, not merely the process exit.\src{status}}

\subsection{Memory store and VA manager}

\doc Guest video memory is one host allocation with offsets corresponding to guest
framebuffer offsets. Guest system memory comes from the VMM's shared backing. A GPU walk
reads guest tables and diffs them against confirmed host placements. Host maps/unmaps apply
the changes, invalidate, and acknowledge runs. A refused map remains a difference.
The committed-placement record is not an adopted guest page table.\src{architecture}
\src{diffcode}\src{memcode}

\attack{Only acknowledged host placements become committed state. Old backing remains owned
until the relevant invalidation acknowledges its retirement.}
{The diff model, apply tests and hardware walk gate compare emitted runs, acknowledgements
and resulting translations. The retirement experiment adds explicit delayed-free tests.
\src{diffcode}\src{applycode}\src{retirementpr}}
{A software model can assume host RM ordering that hardware does not provide. Table roots,
shared subtrees, root withdrawal and failed host invalidation require separate schedules.
Whole-space scanning is not equivalent to remapping every unchanged leaf.\src{latency}}

\subsection{Mirrored scheduling groups}

\readcode Host group/member construction preserves a guest TSG grouping, and twin lifetime follows
the trusted object model.
The late-joiner design documents why a newly added channel cannot silently miss earlier
scheduling state. Address-space mirroring and TSG mirroring are separate concerns: preserving
one does not prove the other.\src{tsgcode}\src{latejoin}\src{kinds}

\attack{Channels in one guest scheduling group retain the corresponding host grouping and
required state as members arrive and leave.}
{Grouping tests and the video-engine validation exercise group creation and related channel
use.\src{tsgcode}\src{video}}
{A test with all channels present at initial configuration misses late arrival. Reused handles
and reset generations can make a correct group lookup refer to the wrong life.\src{latejoin}}

\subsection{User-ring lifecycle: Scan, Pointer and Window}

\doc A registration is a guest request that arms a later software action in its
command stream; that action can move the ring source. Windows user-work can therefore require
software intervention around registration and ring
migration. In Scan mode the worker inspects command packets. Pointer mode forwards admitted
GP references without reading the pushbuffer. Window mode restores scanning across the
registration boundary. An anchor is an observed producer position selected as that boundary.
The clean anchor must have been observed before the registration reply;
a later convenient cursor cannot retroactively establish that boundary.\src{lifecycle}

\attack{Doorbell-covered entries appear once, in order and bit-exact; Pointer reads no
pushbuffer; a validated descriptor is the snapshot actually published.}
{Real-turn model, property and threaded stress suites cover mode transitions, ring migration,
packet continuation, poisoned bytes above PUT and registration at different positions.
\src{lifecycle}\src{turncode}}
{Finite model schedules can omit a real host lock holder, DMA ordering or guest mutation.
A normal-user experiment cannot certify a hostile user that rewrites its own memory.
The boundary remains native USER privilege plus admitted mappings.\src{rulings}}
\begin{figure}[H]
\centering
\begin{tikzpicture}[node distance=20mm]
\node[concept] (scan) {Scan\\inspect packets};
\node[concept,right=of scan] (pointer) {Pointer\\forward references};
\node[concept,right=of pointer] (window) {Window\\inspect registration span};
\draw[flow] (scan) -- node[above=8mm,align=center,font=\scriptsize,text width=32mm]
  {binding known; no registration} (pointer);
\draw[flow] (pointer) -- node[above=8mm,align=center,font=\scriptsize,text width=32mm]
  {registration live at the anchor} (window);
\draw[flow] (window.south) -- ++(0,-15mm) -| (pointer.south);
\node[below=18mm of pointer,align=center,font=\scriptsize,text width=60mm]
  {return only when registration is gone and the relay is quiescent};
\end{tikzpicture}
\caption{Main lifecycle path; direct Scan-to-Window and teardown edges are omitted.
  Registration and packet-boundary rules govern each transition.\src{lifecycle}}
\label{fig:lifecycle}
\end{figure}


\subsection{Lazy observation: a stale view must only delay}

\readcode Pointer mode caches a previously loaded PUT within the same ring generation and
reduces private USERD observations. Scan and Window remain fresh. Low credits, waits, split
segments and a full guest ring force observations. A stale GET may reduce available credit;
it must never increase credit or release a marker gate early. Mode changes and migration
invalidate cached state.\src{lifecycle}\src{lazycode}

\attack{No idle decision relies on cached PUT, and no stale observation is the sole missing
wake for a wait that can already be satisfied.}
{The lazy-turn suites compare forwarding identity with observation enabled and disabled,
exercise small rings, and test forced observations and generation changes.\src{lifecycle}}
{A counter of fewer reads is not proof of liveness. Combined performance arms cannot attribute
speed to each optimisation separately. Hardware completion and holder-release timing remain
outside the pure predicate model.\src{lifecycle}}

\subsection{Relay, handoff and interrupts}

\readcode The relay has explicit progress obligations. If a storage, row or page-table holder
prevents progress, release must arrange another attempt; repeatedly trying a busy lock is not
an event-delivery proof. The baseline includes the relay handoff fix cited by the README and
a later merged lost-wake correction. Experimental split-latency instrumentation keeps that
relay design rather than introducing an attempt limit.\src{readme}\src{turncode}\src{latency}

\attack{Registering a waiter concurrently with completion or lock release cannot lose the
last wake. An interrupt invites a check; it does not manufacture a completed fence.}
{Relay interleaving tests, completion tests and the live regressions cited in the README
exercise parts of this claim. The later probe experiment explicitly diagnoses missing wakes.
\src{readme}\src{probepr}}
{Host events are GPU-wide, so an interrupt can be unrelated to this channel. Observing a marker
late does not timestamp when the GPU wrote it. A worker busy elsewhere can delay observation
without a lost physical event.\src{faq}\src{latency}}

\subsection{Display and the broker}

\doc An emulated display engine accepts the stock driver's display protocol; a separate
broker presents owned frame copies and handles input. The display contract requires guest flips to
complete according to guest semantics while copies select the newest frame independently.
The broker owns cursor composition, including inversion. Its absence or silence must not
freeze guest display progress.\src{architecture}\src{display}\src{rulings}

\attack{Guest flip completion is independent of a stalled viewer; frames and credits remain
bounded; the broker never receives arbitrary guest memory as an owned frame.}
{Display engine tests, plane-output oracles, socket simulations and pacing models exercise
flip order, slow brokers and pixel formats.\src{display}\src{planes}\src{pacing}}
{A fake broker cannot establish compositor, DMA-buffer or colour behaviour on every desktop.
Correct pixels do not prove correct flip latency, and a correct cursor does not prove safe
frame ownership.\src{pacing}\src{security}}
